% ============================================================
\chapter{System Design}
% ============================================================

\section{System Architecture and Program Flow}

NetMind's architecture separates four concerns: deploying the network,
resolving device names to real network identities, deciding what action to
take, and checking that the action is allowed before it runs.

\textbf{Deployment.} Containerlab is used only once per deployment:
\code{generate\_topology.py} builds a topology YAML file from a Python
description of the branches and devices (\code{branches.py}, the single
source of truth for the layout and IP plan), and \code{containerlab deploy}
creates the actual Docker containers and virtual links. \code{make up} runs
the whole sequence, including the idempotent creation of the Linux bridges
that act as switches and hubs. After this point Containerlab is no longer
involved; every device is simply an ordinary Docker container.

\textbf{The device registry.} Because container names and IP addresses are
only known once Containerlab has deployed them, \code{build\_registry.py}
runs \code{containerlab inspect}, cross-references the live state against the
static branch design, and writes a single \code{registry.json} mapping each
device's name and owner to its real container name and IP addresses. This
registry, not Containerlab, is what the rest of the system queries at runtime.

\textbf{Services and backends.} All lab logic lives in three typed services
(\code{DeviceService}, \code{FirewallService} and \code{CameraService} in
\code{services.py}) which use the Docker SDK to start and stop containers
and to execute commands inside them, such as running iptables inside a
firewall or requesting a camera's stream. The outside world (the Docker
client, the registry loader) is passed into each service, so every service
can be tested with plain fakes. The services sit behind a
\code{NetworkBackend} interface with two implementations: a
\code{ContainerlabBackend} for the simulated lab and a
\code{RealLanBackend} for the physical network. Each backend declares the
capabilities it has, and a tool whose capability the backend lacks is
refused in code.

\textbf{The language model layer.} Every tool the model can use is one
registered Python function (\code{tool\_catalog.py}); its schema is built
from the function's name, type hints and docstring, and its decorator records
the capability it needs and its risk level. The web server sends the user's
message and the schemas of the active mode's tools to the model through a
provider interface. The model decides whether to reply directly or to call
one or more tools; each call passes through the safety guard
(Section~\ref{sec:guard}), the result is fed back to the model, and this
repeats for up to five rounds until a final reply is produced. The language
model therefore never touches Docker or Containerlab directly; it only sees
named tools and their results, which is also what lets the same model operate
the separate real-network mode simply by swapping the backend.
Figure~\ref{fig:programflow} shows the deploy-time path and the runtime
loop, and Figure~\ref{fig:layers} shows the layers a tool call passes
through.

\begin{figure}[H]
\centering
\begin{tikzpicture}[
  font=\footnotesize,
  box/.style={draw,rounded corners=2pt,fill=white,align=center,minimum height=1.0cm,inner sep=3pt,text width=2.9cm},
  zl/.style={font=\scriptsize\itshape,anchor=south west,inner sep=2pt},
  note/.style={font=\scriptsize\itshape,text=black!75},
  >={Stealth[length=2mm]}]
  % deploy time
  \draw[dashed] (-0.2,0.9) rectangle (14.4,-0.9);
  \node[zl] at (-0.2,0.9) {Deploy time (\code{make up}, once per deployment)};
  \node[box] (gen) at (1.65,0) {\code{generate\_topology.py}};
  \node[box] (dep) at (5.0,0) {\code{containerlab deploy}};
  \node[box] (reg) at (8.55,0) {\code{build\_registry.py}};
  \node[box,fill=boxgray] (json) at (12.4,0) {\textbf{\code{registry.json}}};
  \draw[->] (gen) -- (dep); \draw[->] (dep) -- (reg); \draw[->] (reg) -- (json);
  % runtime
  \draw[dashed] (-0.2,-1.7) rectangle (14.4,-7.3);
  \node[zl] at (-0.2,-1.7) {Runtime (every chat request)};
  \node[box,fill=userblue] (chat) at (1.6,-3.0) {Browser\\(user message)};
  \node[box] (ws) at (6.3,-3.0) {\textbf{Agent loop}\\(\code{web\_server.py},\\\code{agent\_core.py})};
  \node[box] (llm) at (12.4,-3.0) {\textbf{Groq API}\\tool-calling LLM};
  \node[box,fill=guardorange] (guard) at (6.3,-5.3) {\textbf{Safety guard}\\run, or ask the user\\to confirm; audit};
  \node[box,text width=3.2cm] (be) at (12.4,-5.3) {\textbf{Backend + services}\\lab: Docker SDK\\real LAN: scanner};
  \draw[->] ([yshift=2mm]chat.east) -- node[above,note]{message} ([yshift=2mm]ws.west);
  \draw[->] ([yshift=-2mm]ws.west) -- node[below,note]{reply} ([yshift=-2mm]chat.east);
  \draw[->] ([yshift=2mm]ws.east) -- node[above,note]{messages + tool schemas} ([yshift=2mm]llm.west);
  \draw[->] ([yshift=-2mm]llm.west) -- node[below,note]{tool call or final text} ([yshift=-2mm]ws.east);
  \draw[->] (ws) -- node[left,note]{each tool call} (guard);
  \draw[->] (guard) -- (be);
  \draw[->] (be.south) -- ++(0,-0.7) -| node[pos=0.25,below,note]{result fed back to the model (up to 5 rounds)} (8.35,-3.25) -- (8.35,-3.25 -| ws.east);
  \draw[->,dashed] (json.east) -- (14.05,0) -- node[pos=0.45,sloped,above,note]{device names resolved} (14.05,-5.3) -- (be.east);
\end{tikzpicture}
\caption{NetMind program flow, showing the deploy-time path and the runtime tool-calling loop}
\label{fig:programflow}
\end{figure}

\begin{figure}[H]
\centering
\begin{tikzpicture}[
  font=\footnotesize,
  layer/.style={draw,rounded corners=2pt,text width=13.6cm,minimum height=0.85cm,align=center,fill=boxgray,inner sep=2pt},
  >={Stealth[length=2mm]}, node distance=0.28cm]
  \node[layer,fill=userblue] (ui) {\textbf{Web app} (\code{index.html}): chat, dashboard, devices, alerts, logs, reports};
  \node[layer,below=of ui] (api) {\textbf{REST API} (FastAPI): chat, confirm, devices, alerts, logs, report, \dots};
  \node[layer,below=of api] (agent) {\textbf{Agent loop} (\code{agent\_core.py}), \textbf{LLM provider} (\code{llm.py}), \textbf{sessions} (SQLite)};
  \node[layer,below=of agent,fill=guardorange] (reg) {\textbf{Tool registry} (\code{tools.py}): schema, capability and risk of each tool, checked by the \textbf{safety guard}};
  \node[layer,below=of reg] (be) {\textbf{Network backends} (\code{backends.py}): Containerlab backend (6 capabilities) \;|\; real-LAN backend (inventory only)};
  \node[layer,below=of be] (svc) {\textbf{Services} (\code{services.py}): devices, firewalls, cameras \;|\; \code{network\_scanner.py}};
  \node[layer,below=of svc,fill=white] (world) {Docker containers of the Containerlab lab \;|\; the real local network};
  \foreach \a/\b in {ui/api,api/agent,agent/reg,reg/be,be/svc,svc/world} \draw[->] (\a) -- (\b);
\end{tikzpicture}
\caption{Layers a request passes through after the November refactor}
\label{fig:layers}
\end{figure}

\section{Detailed Design}

\subsection{Detailed Component Description}

\para{Simulated Network Device Control Module}

NetMind's simulated mode operates against a virtual network built with
Containerlab: a core router, three branch firewalls, three switch or hub
segments, eight hosts, and three IP cameras, connected with real IP
addressing and routing between branches. This module exposes five groups of
capabilities to the language model, summarised in Table~\ref{tab:simcaps}.

\textbf{Device status and control.} Reporting whether a host, firewall,
camera, or router is online or offline, retrieving its IP address, and
powering it on or off. Device names are matched exactly by id or owner first
and then, for small spelling mistakes, by a careful fuzzy match that only
accepts one clearly best candidate. Power operations are idempotent, so a
request to power on an already-running device is reported as no change
rather than a false action. Because stopping a container deletes the virtual
link Containerlab created, powering a host or camera back on re-attaches it
to its switch and replays its addressing and routes, then checks that it can
reach its gateway.

\textbf{Network visibility and connectivity.} Listing every device with its
owner, type, and status, listing the devices on a specific switch or hub
segment, validating the segment name against the real topology rather than
silently returning an empty result, and pinging one device from another over
the lab's data network (not the management network, which would hide a
broken link).

\textbf{Firewall auditing.} Each branch firewall runs a real Linux iptables
firewall rather than a placeholder. NetMind can list its configured rules,
report open ports, audit it for common misconfigurations such as an
administrative port left open to any source, and block or allow a port on
request; changes persist across a normal power cycle. Rules are applied to
both the INPUT chain and the FORWARD chain, because a rule on INPUT alone
only protects the firewall itself and does not stop traffic passing through
it between branches; an integration test confirms this with real traffic
(Section~\ref{sec:testing}). Rule changes are idempotent: NetMind checks
whether a rule already exists before inserting it, so repeating a request
never creates duplicates.

\textbf{Camera security.} Each camera runs a small real HTTP service, so
NetMind can check whether its stream is actually reachable independently of
the device's power state, start or stop the stream service, and check or
change the camera's default password.

{\small
\begin{xltabular}{\textwidth}{|L{3.6cm}|Y|}
\hline
\head \textbf{Capability Group} & \textbf{What It Provides} \\ \hline
\endhead
Device status and control & Online/offline status, IP address, and idempotent power on/off for any host, firewall, camera, or router, with automatic re-attachment after power-on \\ \hline
Network visibility & Listing every device, or every device on a specific switch/hub segment, validated against the real topology \\ \hline
Connectivity & Pinging one lab device from another over the data network \\ \hline
Firewall auditing & Listing configured iptables rules (INPUT and FORWARD chains), reporting open ports, auditing for misconfigurations, and blocking/allowing ports \\ \hline
Camera security & Checking stream reachability independently of power state, starting/stopping the stream, and checking/changing the default password \\ \hline
\caption{Simulated-mode capability groups (summary)}\label{tab:simcaps}\\
\end{xltabular}}

All of this runs over genuine IP addressing and routing between branches
rather than a flat management network, verified through packet round-trip
testing and Time-To-Live analysis to confirm that traffic between branches
actually traverses the intended switch, firewall, and router path.

\para{Real Network Discovery Module}

The real network discovery module is the component that gives NetMind
visibility of an actual physical network. It is written in Python and is
deliberately read-only: it observes the local network and reports what it
finds, but it never changes the state of any device. This is what separates
it from the device control module, which manages the Docker containers of the
simulated Containerlab network. Keeping the two apart means that every
operation which can change a device is confined to the simulated
environment, while the real network can only be observed, which removes any
risk of the language model accidentally altering a live network. The
boundary is enforced by the backend: \code{RealLanBackend} has only the
inventory capability, so the power, firewall and camera tools cannot run
against it whatever the model asks.

The module exposes three functions to the tool layer: listing the devices on
the network, checking whether a single device is online, and resolving a
hostname to an IP address. Each device is returned once, as a dictionary
containing its IPv4 address, its IPv6 addresses, its MAC address, its
hostname, its status, and \code{found\_by}, the list of discovery methods
that saw it. The dashboard shows \code{found\_by} for every device.

\textbf{Program flow.} The first version of the module ran ARP and fell back
to a ping sweep only when ARP could not run, so a scan reported the devices
found by one method only, even though each method finds devices the other
misses. In November the scan was redesigned so that all methods run at the
same time and their results are merged. Table~\ref{tab:scansteps} lists the
stages of a scan and Figure~\ref{fig:scanflow} shows the flow.

{\small
\begin{xltabular}{\textwidth}{|>{\centering\arraybackslash}p{0.8cm}|L{3.2cm}|Y|}
\hline
\head \textbf{Step} & \textbf{Stage} & \textbf{Description} \\ \hline
\endhead
1 & Subnet detection & The local subnet is derived from the machine's own IP address on its default route, so that no configuration is needed from the user. \\ \hline
2 & Parallel discovery & Four methods start together: a broadcast ARP request to every address in the subnet (IP and MAC of each device that answers); a ping sweep of the subnet with 64 parallel workers; an mDNS browse of six seconds that maps addresses to the names devices advertise; and a ping to the IPv6 all-nodes address so that IPv6 neighbours appear in the kernel's neighbour table. A method that cannot run, such as ARP without raw-socket permission, is reported as not run, and the others still complete. \\ \hline
3 & Merge & Devices are merged by IPv4 address. Each keeps the list of methods that found it. A device found only by ping takes its MAC address from the kernel's neighbour table. \\ \hline
4 & IPv6 matching & IPv6 addresses from the neighbour table are attached to the device with the same MAC address. \\ \hline
5 & Hostname resolution & For each device the hostname is taken from the mDNS table if present, and otherwise from a reverse DNS lookup, done in parallel. \\ \hline
6 & Output & The list is sorted by address and returned with the methods that ran; it is cached for 20 seconds. \\ \hline
\caption{Real-network discovery module program flow steps}\label{tab:scansteps}\\
\end{xltabular}}

\begin{figure}[H]
\centering
\begin{tikzpicture}[
  font=\footnotesize,
  box/.style={draw,rounded corners=2pt,fill=white,align=center,minimum height=0.85cm,inner sep=3pt},
  par/.style={box,fill=boxgray,text width=2.6cm},
  >={Stealth[length=2mm]}]
  \node[box,text width=7.4cm] (sub) at (6.3,0) {1. Subnet detection (own IP on the default route)};
  \node[par] (arp)  at (1.5,-2.1) {ARP broadcast\\(IP + MAC)};
  \node[par] (ping) at (4.7,-2.1) {Ping sweep\\(64 workers)};
  \node[par] (mdns) at (7.9,-2.1) {mDNS browse\\(6 seconds)};
  \node[par] (ndp)  at (11.1,-2.1) {IPv6 all-nodes\\ping};
  \draw[dashed] (-0.1,-1.25) rectangle (12.7,-2.85);
  \node[font=\scriptsize\itshape,anchor=north west,inner sep=2pt] at (-0.1,-2.85) {2. Parallel discovery (a method that cannot run is skipped)};
  \node[box,text width=8.4cm] (merge) at (6.3,-4.1) {3. Merge by IPv4 address and record \code{found\_by};\\MAC for ping-only hosts from the kernel neighbour table};
  \node[box,text width=8.4cm] (v6) at (6.3,-5.45) {4. Attach the IPv6 addresses that have the same MAC};
  \node[box,text width=8.4cm] (host) at (6.3,-6.55) {5. Hostname: mDNS table, then reverse DNS};
  \node[box,text width=8.4cm,fill=guardorange] (out) at (6.3,-7.65) {6. Output: one entry per device, sorted, cached for 20 s};
  \foreach \n in {arp,ping,mdns,ndp} \draw[->] (sub.south) -- (\n.north);
  \draw[->] (8.5,-2.85) -- (8.5,-3.45 |- merge.north);
  \draw[->] (merge) -- (v6); \draw[->] (v6) -- (host); \draw[->] (host) -- (out);
\end{tikzpicture}
\caption{Real-network discovery module program flow}
\label{fig:scanflow}
\end{figure}

\textbf{Design Decisions and Alternatives Considered.} This part explains why
each technique and library was chosen, and which alternatives were considered
and not selected (summary in Table~\ref{tab:decisions}).

{\small
\begin{xltabular}{\textwidth}{|L{3.2cm}|Y|}
\hline
\head \textbf{Decision} & \textbf{Chosen Approach (Alternatives Considered)} \\ \hline
\endhead
Subnet detection & UDP socket trick, standard library only, with \code{ip addr} as a fallback (netifaces/psutil: accurate but an extra dependency; gethostbyname: often returns 127.0.0.1) \\ \hline
Host discovery & ARP via Scapy and a ping sweep, run together and merged (Nmap: more capable but a heavier external binary; arp-scan: external dependency; OS ARP cache alone: misses undiscovered hosts; passive sniffing: slow) \\ \hline
IPv6 discovery & Ping the all-nodes multicast address, then read the kernel neighbour table (a full IPv6 sweep: impossible, a /64 has $2^{64}$ addresses) \\ \hline
Concurrency & ThreadPoolExecutor: one pool for the four methods, 64 workers for the ping sweep (multiprocessing: heavier for I/O-bound work; asyncio: would need a restructured codebase) \\ \hline
Hostname discovery & mDNS via zeroconf (NetBIOS: Windows-only; SSDP/UPnP: fewer device types; DHCP leases: needs router access; SNMP: needs credentials) \\ \hline
Reverse DNS & Python standard library, used as a secondary source \\ \hline
Status checking & A single ping (TCP connect or ARP probe identified as possible improvements) \\ \hline
Optional dependencies & Scapy and zeroconf imported in try/except, so the module degrades gracefully if either is missing \\ \hline
Privileges & A copy of the virtualenv's Python with only the raw-socket capability, so ARP works without running the server as root \\ \hline
\caption{Design decisions and alternatives considered (summary)}\label{tab:decisions}\\
\end{xltabular}}

\textbf{Subnet detection.} The active interface is found by opening a UDP
socket towards a public address and reading the local address the operating
system assigns to it. No packet is actually sent, and the method relies only
on the Python standard library. If there is no default route at all, the
module reads the interface configuration with \code{ip addr} instead.
Alternatives were the netifaces and psutil libraries, which are more accurate
on machines with several interfaces but add an external dependency, and
gethostbyname applied to the machine's own hostname, which frequently returns
the loopback address 127.0.0.1 and is therefore unreliable. The chosen method
assumes a /24 subnet, which is a known limitation.

\textbf{Host discovery with ARP, Scapy and a ping sweep.} ARP was chosen
because it works at the data link layer: a device cannot ignore an ARP
request and still take part in an Ethernet network, so ARP finds devices that
are configured to block ICMP ping, and it also returns the MAC address. Scapy
was chosen to send the requests because it is a pure Python library that
integrates directly with the rest of the code and gives full control over
packet construction without an external program. In practice, however, some
devices answer ping but miss the short ARP window, so the ping sweep is not
only a fallback: running both and merging the results finds more devices than
either method alone, and \code{found\_by} makes visible which method saw each
device. The alternatives were Nmap, used directly or through python-nmap,
which is far more capable but is a separate binary larger than this module
requires; the arp-scan command-line tool, which is fast but is again an
external dependency; the operating system's ARP cache, which only lists
devices that the machine has already communicated with; and passive sniffing
of network traffic, which needs no probes but can take a long time to reveal
quiet devices.

\textbf{IPv6.} An IPv6 subnet is far too large to sweep, so the module sends
one ping to the IPv6 all-nodes multicast address on the active interface. Every
IPv6 device on the link answers, which places it in the kernel's neighbour
table; the module then reads that table and attaches each IPv6 address to the
device with the same MAC address. A device often has several IPv6 addresses
(a link-local \code{fe80::} address and one or more global ones), and all are
listed.

\textbf{Concurrency.} Checking 254 addresses one after another would take
several minutes, since each unanswered ping waits up to one second. A
ThreadPoolExecutor with 64 workers runs the pings in parallel and reduces the
sweep to a few seconds, and a second small pool runs the four discovery
methods at the same time, so a full scan takes about as long as its slowest
method. Threads suit this task because it is bound by waiting on the network
rather than by computation. The alternatives were multiprocessing, whose
separate processes add memory and start-up cost for work that is only
waiting, and asyncio, which would be efficient but would require the
surrounding code to be restructured around asynchronous functions.

\textbf{Hostname discovery with mDNS.} Multicast DNS was chosen because small
and home networks often have no reverse DNS records, so a reverse lookup
alone returns no name for most devices. mDNS lets printers, AirPlay devices,
Chromecast devices, and shared workstations announce their names without any
credentials or access to the router, and the zeroconf library provides a pure
Python implementation. The browse runs on every scan rather than once at
server start, so that a device that joins the network later is still found.
The alternatives were NetBIOS name queries, which are mostly limited to
Windows hosts; SSDP and UPnP discovery, which cover fewer device types; DHCP
lease tables read from the router, which need administrative access to the
router; and SNMP queries, which need SNMP to be enabled and credentials to be
known. A limitation of mDNS is that a device which advertises no service, such
as an idle phone, appears without a hostname, although it is still listed by
IP and MAC address.

\textbf{Reverse DNS and status checking.} Reverse DNS through the standard
library is used as a secondary source, because it is already available in
Python at no cost and gives names on networks that do maintain DNS records.
The status of a single device is determined with one ping, which keeps it
consistent with the ping sweep and needs no extra dependency. A TCP
connection check or an ARP probe would be more reliable against devices that
block ICMP, and are identified as possible improvements.

\textbf{Optional dependencies and output format.} Scapy and zeroconf are
imported inside try and except blocks, and their availability is recorded in
flags. If either library is missing the module degrades gracefully instead of
failing to start, which makes it easier to run on different machines. Results
are returned as plain dictionaries, rather than custom classes, because they
map directly to JSON for the API.

\textbf{Limitations.} ARP scanning only works within the local broadcast
domain and needs the raw-socket capability. The ping command options used are
those of Linux and would have to be adapted for Windows and macOS. A device
that blocks ICMP may be reported as offline by the status check. The subnet is
assumed to be a /24. Under WSL2's default network mode the scanner sees the
virtual network between Windows and WSL rather than the physical LAN, so real
mode needs WSL's mirrored networking mode or a native Linux machine. Finally,
network scanning should only be carried out on networks owned by the user or
for which permission has been given.

\para{Live Monitoring Dashboard Module}

The live monitoring dashboard is the module that gives the operator a single,
continuously updated view of the network without having to ask the chat
interface a question first. It is implemented as part of a unified
single-page application that also hosts the chat interface, so both views
share one navigation shell, one theme, one language setting, and one REST
client rather than being two separate pages. Table~\ref{tab:dashcaps}
summarises its capability groups.

{\small
\begin{xltabular}{\textwidth}{|L{3.6cm}|Y|}
\hline
\head \textbf{Capability Group} & \textbf{What It Provides} \\ \hline
\endhead
Unified application shell & One entry point providing both the chat and dashboard views, with shared sidebar navigation, light/dark theme, Arabic/English language switching, and a single REST client used by both views \\ \hline
Mode-aware refresh & Polls \code{/api/devices} and \code{/api/alerts} (simulated) or \code{/api/real-devices} and \code{/api/real-alerts} (real network) on a fixed interval, discarding any response that no longer matches the currently selected mode to avoid a stale-data race \\ \hline
Network visibility & Renders device and connection state, live statistics (online/offline counts), a topology view, and an active-alerts panel from the same registry data used by the chat tools \\ \hline
Device inventory view & A dedicated devices page listing currently active devices, with best-effort device-type classification (laptop, phone, router, etc.) from MAC OUI prefixes and hostname heuristics, each shown with a type icon; in real mode it also shows how each device was found \\ \hline
Evidence chain & Each chat reply that triggers one or more backend tool calls is accompanied by a collapsible, step-by-step trace of which tool was called, with what arguments, and what it returned, so the operator can inspect the reasoning behind a diagnosis rather than trusting it blindly \\ \hline
\caption{Live monitoring dashboard capability groups (summary)}\label{tab:dashcaps}\\
\end{xltabular}}

\textbf{Design decisions.} The dashboard and chat interface were originally
built as two separate HTML pages and were later merged into one application
shell, each view kept in its own section and toggled client-side without a
page reload, so that navigation state, theme, language, and the active
network mode persist across both views instead of resetting. The refresh
logic is deliberately race-safe: each poll captures which mode it was
requested for, and a response is applied to the interface only if that mode
is still the one selected when the response arrives, which prevents
simulated-network data from briefly overwriting real-network data (or vice
versa) when the operator switches modes while a request is in flight.
Device-type classification in the inventory view is explicitly best-effort
only: modern phones that randomise their MAC address (iOS 14+/Android 10+)
cannot be reliably identified by OUI lookup, so the interface documents this
limitation rather than presenting the classification as certain.

\textbf{Limitation.} Like the rest of the frontend, the dashboard has no
authentication and assumes a single trusted operator, consistent with the
scoped-out security requirement (NFR12) already noted in Chapter~3.

% ------------------------------------------------------------
\subsection{Safety Guard: Confirmation, Protected Devices and Audit Log}
\label{sec:guard}

A language model that can call tools can also be wrong or be misled by a
request, so NetMind does not let the model carry out destructive actions
itself. Every tool is registered with one of three risk levels. \emph{Read}
tools (status, listings, audits, pings) run immediately. \emph{Action} tools,
which change state but are easy to undo (power on, start or stop a camera
stream), also run immediately and are audited. \emph{Destructive} tools
(power off, block port, allow port, and change camera credentials) never run
on the model's call. Because the risk level is part of the tool's
registration, the guard, the backends and the dashboard all read it from the
same table, and a new tool cannot be added without one.

When the model calls a destructive tool, the guard does not run the action.
It stores a pending action containing the tool, its target, and its
arguments, and returns a ``pending'' result to the model. The pending action
is shown to the user in the chat as a Confirm/Cancel box. The action runs only
when the user presses Confirm, which sends a request to a separate endpoint
(\code{POST /api/confirm}) that the model has no way to call. A pending action
that is not answered within five minutes expires.
Figure~\ref{fig:confirm} shows the sequence.

The core router and every firewall are classed as protected devices.
Confirming an action on a protected device requires the user to tick an
additional acknowledgement box, and the server rejects the request without
it.

Every outcome (executed, cancelled, expired, or rejected) is written to an
audit log, stored in the SQLite database and also appended to
\code{logs/audit.jsonl}. Each entry records the time, the tool, the target,
the outcome, and the result. Camera passwords are redacted before an entry is
written.

Language models sometimes answer an action request in text, for example
saying that a request is ``pending'', without calling any tool. In that case
the user would see no Confirm box and nothing would happen. To catch this,
the agent loop checks the reply against two patterns: a destructive request
that produced no tool call, and a reply that claims something is pending or
done. When either matches, the loop retries once with an explicit instruction
to call the tool. This retry is a safety net for usability and does not weaken
the guard, because no path from the model leads to execution.

\begin{figure}[H]
\centering
\begin{tikzpicture}[
  font=\small,
  stp/.style={draw,rounded corners=2pt,align=center,minimum width=6.2cm,minimum height=0.85cm,inner sep=3pt},
  user/.style={stp,fill=userblue},
  model/.style={stp,fill=boxgray},
  guard/.style={stp,fill=guardorange},
  side/.style={font=\scriptsize\itshape,align=left,anchor=west},
  >={Stealth[length=2mm]}, node distance=0.42cm]
  \node[user] (a) {User asks in the chat:\\``turn off host-khalid''};
  \node[model,below=of a] (b) {Language model calls\\the \code{power\_off} tool};
  \node[guard,below=of b] (c) {Guard stores a \textbf{pending action} and returns\\``pending'' to the model (nothing has run)};
  \node[user,below=of c] (d) {Chat shows a \textbf{Confirm / Cancel} box\\(protected devices: extra acknowledgement)};
  \node[user,below=of d] (e) {User presses \textbf{Confirm}};
  \node[guard,below=of e] (f) {Guard runs the tool and records the outcome};
  \node[stp,below=of f] (g) {Audit log entry (SQLite + \code{logs/audit.jsonl})};
  \foreach \x/\y in {a/b,b/c,c/d,d/e,e/f,f/g} \draw[->] (\x) -- (\y);
  \node[side] at ($(c.east)+(0.3,0)$) {expires after 5 minutes\\if unanswered};
  \node[side] at ($(e.east)+(0.3,0)$) {\code{POST /api/confirm}, a route\\the model cannot call};
\end{tikzpicture}
\caption{Confirmation flow for destructive actions: the model can only create a pending action; only the user's button press executes it}
\label{fig:confirm}
\end{figure}

% ------------------------------------------------------------
\subsection{Tool Registry, Network Backends and Model Provider}

In the first version, the model-facing description of every tool was written
by hand in the web server, and separate tables mapped tool names to functions
for each mode. The same information therefore lived in several places, and
the rule that real mode is read-only was only stated in the prompt. Following
the supervisor's assessment, this part of the system was refactored in
November into three small, independently tested pieces.

\textbf{Tool registry.} A tool is one plain Python function decorated with
the capability it needs and its risk level. The schema the model receives is
generated from the function itself: its name, its type hints, and the
parameter descriptions in its docstring. Adding a tool means writing one
function; nothing else lists tools. The simulated mode currently offers
seventeen tools, and the real mode three.

\textbf{Network backends.} The \code{NetworkBackend} interface describes a
network NetMind can talk to. \code{ContainerlabBackend} declares six
capabilities (inventory, segments, diagnostics, power, firewall and camera);
\code{RealLanBackend} declares only inventory. The set of tools offered to the
model in each mode is derived from the backend's capabilities, and calling a
tool whose capability the backend lacks raises an error in code. The
read-only guarantee for the real network (NFR4) therefore no longer depends
on the prompt. The same interface is also the extension point for future
backends, such as virtual machines or real network devices.

\textbf{Model provider.} The agent loop only calls one method,
\code{complete(messages, tools)}, on an \code{LLMProvider}. The Groq provider
implements it with an OpenAI-compatible client, so changing the model vendor
means writing one small class, which is what the model comparison planned for
Project 2 relies on. The Groq client and all settings are created on first
use rather than when the program is imported, so the application, and its
tests, start even without an API key or Docker.

% ------------------------------------------------------------
\subsection{Persistence and Least Privilege}

\textbf{Persistence.} Chat sessions, their messages, and the audit log are
stored in a SQLite database (\code{logs/netmind.db}), created readable by its
owner only because it contains chat text. One connection is shared by the
server's threads and guarded by a lock. A storage failure never breaks a
request: the agent catches it and carries on in memory. The alert history is
kept in \code{logs/alerts.json}, so a device that is down across a restart
keeps its original alert and its downtime keeps counting.

\textbf{Running without root.} The server needs exactly two privileges: a raw
socket for ARP scanning, and root rights to re-attach a lab device to its
bridge after power-on. Rather than running the whole server with
\code{sudo}, \code{make setup-perms} grants each one narrowly. A private copy
of the virtualenv's Python receives only the raw-socket capability
(\code{CAP\_NET\_RAW}), and \code{make run} uses it when present. Re-attaching
is done by one small, root-owned helper script, \code{netmind-reconnect},
which the user may run through \code{sudo} without a password and which
checks every argument (the container must belong to the lab, the bridge must
exist, and names may contain only safe characters) before doing anything.

% ------------------------------------------------------------
\subsection{Web API, Alerts and Notifications}

The backend exposes a small REST API, listed in Appendix~\ref{app:api}. The
chat endpoint runs the agent loop, bounded at five tool-calling rounds. The
dashboard uses read endpoints for devices, alerts, and logs, and a power
endpoint for the on/off buttons (turning on runs directly; turning off returns
a pending action that goes through the guard).

The first version of the dashboard polled the Docker daemon device by device,
which made one refresh take more than fifteen seconds with the full lab
running. Four changes made refreshes fast enough for regular polling:

\begin{enumerate}
\item all container states are read from a single container listing;
  per-container statistics, the slowest call and the only source of the CPU
  figures, are no longer requested, and the dashboard shows status rather than
  CPU usage;
\item the computed result is cached for three seconds, and concurrent callers
  share one computation instead of each starting their own;
\item firewall audits and camera credential checks used for the alert list
  are cached for ten seconds, and the cache is cleared whenever a rule,
  password or power state changes;
\item the frontend coalesces overlapping refreshes and releases a pressed
  button when its action finishes, so a button can no longer remain stuck.
\end{enumerate}

Three views sit on top of this state. The \emph{alert history} records a
``Device Offline'' alert when a device goes down and a ``Device Back Online''
alert, with the measured downtime, when it returns; the earlier alert is
closed but never deleted, so the history shows both that a device went down
and how long recovery took. Firewall findings and cameras still using their
default password raise alerts the same way, each stamped once with the time
it was first seen. The \emph{activity log} is built from the audit log and
lists the actions people and the assistant requested (power changes, port
changes, camera password changes), who requested them and how each one
ended: awaiting confirmation, confirmed, cancelled, denied or expired. A
\emph{report} view summarises open security findings, outages, average and
longest recovery time and operator activity, and can be printed or saved as a
PDF.

When a Telegram bot is configured, every alert and every recovery is also
sent as a message containing the network, the device, the event, its severity
and the time, and real-network devices joining or leaving the network are
reported the same way. A test button in the interface sends a test message,
so the setup can be checked in one click. The integration is optional: without
the two settings it does nothing, and a failure to send never breaks the
request that triggered it.

The web server accepts requests only from its own origin and listens on the
local machine only. A request carrying a different \code{Origin} header is
refused, so a web page open in the same browser cannot power a device off or
confirm a pending action through the user's session.

% ------------------------------------------------------------
\section{Verification and Testing}
\label{sec:testing}

Verification used four layers, summarised in Table~\ref{tab:verification}.

{\small
\begin{xltabular}{\textwidth}{|L{3.6cm}|Y|L{2.4cm}|}
\hline
\head \textbf{Layer} & \textbf{What it checks} & \textbf{Result} \\ \hline
\endhead
Unit and endpoint tests (pytest, 153 test functions in 17 files) & Argument parsing, session pruning and history trimming, the guard (pending, confirm, cancel, expiry, protected acknowledgement, password redaction), the retry safety net, the tool registry and generated schemas, backend capabilities, the services with a fake Docker layer, firewall rule parsing and rule application on both chains, SQLite storage, settings, the permission helper's argument checks, the merged real-network scan, API endpoints, and dashboard caching & All passing in CI, with test coverage of at least 70\% enforced \\ \hline
Browser UI tests (Playwright, 3 tests) & The real web server with a fake Docker layer, driven by a headless browser: powering a device off without its status flashing back online, the alert history, logs and report after an outage, and the chat view & Passing in CI \\ \hline
Integration test against the live lab (1 test) & A block rule on a firewall really stops traffic between branches, and the allow rule restores it & Passed on the team's lab \\ \hline
Manual browser verification & Pending, Confirm and Cancel flow, protected-device acknowledgement, audit log entries, power buttons after off/on, Arabic view, Real Network mode device list with \code{found\_by} & Passed \\ \hline
\caption{Verification layers and results}\label{tab:verification}\\
\end{xltabular}}

Static checks (ruff) and the pytest suite, with a minimum coverage of 70\%,
run automatically in GitHub Actions on every pull request. A dedicated test checks that the unit
tests cannot reach a Docker daemon at all, which proves that they run anywhere
without the lab or an API key. The integration test needs the
lab and is run on demand.
